\section{Problematización}

La investigación aborda la distancia entre los recursos monetarios de los hogares y el costo local de satisfacer necesidades básicas en Chile, Brasil y Argentina. Aunque los tres países cuentan con encuestas de hogares y estadísticas oficiales de precios, las fuentes difieren en los períodos observados, los conceptos de ingreso, la cobertura territorial y la construcción de sus canastas. En consecuencia, no es metodológicamente suficiente yuxtaponer montos nominales o porcentajes nacionales de pobreza para afirmar que el bienestar es directamente comparable.

El problema empírico consiste en construir una medida transparente de cobertura: relacionar el ingreso de los hogares con el costo mensual de una canasta o referencia de consumo para períodos y territorios compatibles. La disponibilidad de microdatos y documentación es desigual. En particular, los archivos de gasto de la EPF chilena describen patrones observados de consumo y no equivalen automáticamente a una canasta normativa; la cesta DIEESE de Brasil es una referencia independiente por capital, mientras que la CBA/CBT de Argentina se publica para el Gran Buenos Aires. La cobertura territorial y las unidades deberán compatibilizarse antes de estimar resultados.

La relevancia comparativa radica en analizar cómo se expresa la capacidad de cubrir necesidades básicas bajo distintos contextos laborales, distributivos e institucionales \parencite{GindlingRonconi2025, Stampinietal2025, cepalPanoramaSocialAmerica2025}. La contribución metodológica esperada es documentar las fuentes, mantener visibles las decisiones de armonización y distinguir las estimaciones descriptivas nacionales de las comparaciones que los datos permiten sostener.